% Shared pleading-paper preamble for CGC-25-631802 TAC-RESPONSE court PDFs.
\documentclass{article}
\usepackage[utf8]{inputenc}
\usepackage{graphicx}
\usepackage{geometry}
    \geometry{top=1in, bottom=2.2in, left=1.0in, right=1.0in, textheight=336pt}
    \setlength{\footskip}{55pt}
\usepackage{ulem}
\usepackage{fancyhdr}
    \renewcommand{\headrulewidth}{0pt}
    \renewcommand{\footrulewidth}{0pt}
\usepackage{parskip}
    \setlength{\parskip}{0mm}
\usepackage{quoting}
    \quotingsetup{leftmargin=0.5in, rightmargin=0.5in, vskip=-1.5mm}
    \makeatletter
        \g@addto@macro\quoting\singlespacing
        \g@addto@macro\quoting{\vspace{-2mm}}
    \makeatother
\usepackage{mathptmx}
\usepackage{amssymb}
\renewcommand{\rmdefault}{ptm}
\renewcommand{\normalsize}{\fontsize{12}{12}\selectfont}
\newcommand*{\ptm}{\fontfamily{ptm} \selectfont \fontsize{12}{12}\selectfont}
\usepackage{setspace}
    \doublespacing
\newcommand{\tab}{\hspace*{.0in}}
\newenvironment{tightcenter}{%
    \setlength\topsep{0pt}
    \setlength\parskip{0pt}
    \begin{center}
}{%
    \end{center}
}
\usepackage{tikz}
\usetikzlibrary{calc}
\AddToHook{shipout/background}{%
    \begin{tikzpicture}[remember picture, overlay]
        \draw[line width=0.5pt] ($(current page.north west) + (0.75in,0)$) -- ($(current page.south west) + (0.75in,0)$);
        \draw[line width=0.5pt] ($(current page.north west) + (0.80in,0)$) -- ($(current page.south west) + (0.80in,0)$);
    \end{tikzpicture}%
}
\usepackage[left, pagewise]{lineno}
\setlength{\linenumbersep}{0.35in}
\renewcommand{\linenumberfont}{\normalfont\small}
\usepackage{titlesec}
\usepackage[hidelinks]{hyperref}
\usepackage{bookmark}
\usepackage{textcomp}
\usepackage{longtable}
\usepackage{booktabs}
\usepackage{array}
\usepackage{multirow}
\usepackage{calc}
% Pandoc pipe-table column widths use \real{#}; map to pgf math (pgf loaded via tikz).
\newcommand{\real}[1]{\pgfmathparse{#1}\pgfmathresult}
\titleformat{\section}{\fontsize{15}{15}\selectfont\normalfont\bfseries\uppercase}{}{0em}{}
\titleformat{\subsection}{\fontsize{15}{15}\selectfont\normalfont\bfseries\uppercase}{}{0em}{}
\titleformat{\subsubsection}{\normalfont\bfseries}{}{0em}{}
\titlespacing*{\section}{0pt}{0pt}{0pt}
\titlespacing*{\subsection}{0pt}{0pt}{0pt}
\titlespacing*{\subsubsection}{0pt}{0pt}{0pt}
\newcommand{\calrules}[1]{Cal. Rules of Court, rule #1}
\newcommand{\ccp}[1]{Code Civ. Proc., \textsection~#1}
\newcommand{\evidcode}[1]{Evid. Code, \textsection~#1}
\newcommand{\vehcode}[1]{Veh. Code, \textsection~#1}
\providecommand{\tightlist}{%
  \setlength{\itemsep}{0pt}\setlength{\parskip}{0pt}}
% --- TOC / TOA hyperlink infrastructure ---
\newcommand{\tacctoclink}[2]{\hyperref[toc:#1]{#2}}
% In-text citation link to authority label (target is usually TABLE OF AUTHORITIES row).
\newcommand{\citecase}[2]{\textit{\hyperref[auth:#1]{#2}}}
\newcommand{\citestatute}[2]{\hyperref[auth:#1]{#2}}
\newcommand{\toclink}[2]{\hyperref[toc:#1]{#2}}
% Authority entries in TABLE OF AUTHORITIES (#1 label slug, #2 italic case name or title, #3 citation tail)
\newcommand{\caseentry}[3]{\par\noindent\hangindent=2em\hangafter=1%
  \phantomsection\label{auth:#1}\textit{#2}, #3\par}
\newcommand{\statuteentry}[3]{\par\noindent\hangindent=2em\hangafter=1%
  \phantomsection\label{auth:#1}#2\ #3\par}
\newcommand{\ruleentry}[2]{\par\noindent\hangindent=2em\hangafter=1%
  \phantomsection\label{auth:#1}#2\par}
\makeatletter
\newlength{\CourtTOCmarkwd}
\newlength{\CourtTOCpagewd}
\newlength{\CourtTOCtitlewd}
\setlength{\CourtTOCmarkwd}{2.8em}
\setlength{\CourtTOCpagewd}{4.2em}
\newcommand{\CourtTOCRow}[3]{%
  \begingroup
  \setlength{\CourtTOCtitlewd}{\dimexpr\linewidth - \CourtTOCmarkwd - \CourtTOCpagewd - 0.6em\relax}%
  \par\nobreak\noindent
  \begin{minipage}[b]{\CourtTOCmarkwd}\raggedright\strut #1\strut\end{minipage}\hspace{0.3em}%
  \begin{minipage}[b]{\CourtTOCtitlewd}\raggedright\strut #2\strut\end{minipage}\hspace{0.3em}%
  \begin{minipage}[b]{\CourtTOCpagewd}\raggedleft\strut #3\strut\end{minipage}\par
  \endgroup
  \vspace{1pt}%
}
\newcommand{\CourtTOCSubitem}[2]{\CourtTOCRow{}{#1}{#2}}
\newcommand{\CourtTOCSubsubitem}[2]{\CourtTOCRow{}{\quad\quad #1}{#2}}
\makeatother

\begin{document}
\nolinenumbers
\begin{singlespace}
Franciscus Dylan Rosario, \\
7624 Melody \\
Rohnert Park, CA 94928 \\
E: soltrinox@gmail.com \\
Plaintiff, In Pro Per \\
\end{singlespace}
\vspace*{9mm}
\begin{tightcenter}
\textbf{SUPERIOR COURT OF THE STATE OF CALIFORNIA} \\
\textbf{COUNTY OF SAN FRANCISCO}
\end{tightcenter}
\vspace*{2.25mm}
\nolinenumbers
\begin{minipage}[t]{3in}
    \raggedright
    \textbf{Franciscus Dylan Rosario,} \\ \textbf{PLAINTIFF}, \\
    \hspace{1cm} v. \\
    \small
    CSAA Insurance Exchange; \\
    and DOES 1 through 50, inclusive; \\
    \normalsize
    \textbf{DEFENDANTS}
    \makebox[3in]{\hrulefill}
\end{minipage}%
\hspace{3mm}
\begin{minipage}[t]{0.5pt}
    \vspace{0pt}
    \rule{0.5pt}{2.0in}
\end{minipage}
\hspace{3mm}
\begin{minipage}[t]{3in}
    \raggedright
    \noindent \textbf{Case No.: CGC-25-631802} \\
    ~\\
    \textbf{MEMORANDUM OF POINTS AND AUTHORITIES} \\
    \textbf{IN SUPPORT OF LEAVE TO FILE THIRD AMENDED COMPLAINT}
\end{minipage}
\vspace*{6mm}
\linenumbers
\begin{center}\textbf{TABLE OF CONTENTS}\end{center}
\vspace{2mm}
\CourtTOCRow{I.}{\tacctoclink{i-introduction}{I. INTRODUCTION}}{\pageref{toc:i-introduction}}
\CourtTOCRow{II.}{\tacctoclink{ii-statement-of-facts}{II. STATEMENT OF FACTS}}{\pageref{toc:ii-statement-of-facts}}
\CourtTOCRow{III.}{\tacctoclink{iii-legal-standard}{III. LEGAL STANDARD}}{\pageref{toc:iii-legal-standard}}
\CourtTOCRow{IV.}{\tacctoclink{iv-the-tac-cures-all-seven-defense-objections}{IV. THE TAC CURES ALL SEVEN DEFENSE OBJECTIONS}}{\pageref{toc:iv-the-tac-cures-all-seven-defense-objections}}
\CourtTOCRow{V.}{\tacctoclink{v-specific-answers-to-defenses-legal-arguments}{V. SPECIFIC ANSWERS TO DEFENSE'S LEGAL ARGUMENTS}}{\pageref{toc:v-specific-answers-to-defenses-legal-arguments}}
\CourtTOCRow{VI.}{\tacctoclink{vi-no-prejudice-to-defendant}{VI. NO PREJUDICE TO DEFENDANT}}{\pageref{toc:vi-no-prejudice-to-defendant}}
\CourtTOCRow{VII.}{\tacctoclink{vii-conclusion}{VII. CONCLUSION}}{\pageref{toc:vii-conclusion}}
\vspace{4mm}

\begin{center}\textbf{TABLE OF AUTHORITIES}\end{center}
\vspace{2mm}
\subsubsection*{Cases}
\caseentry{howard}{Howard v. County of San Diego}{(2010) 184 Cal.App.4th 1422}
\caseentry{kolani}{Kolani v. Gluska}{(1998) 64 Cal.App.4th 402}
\caseentry{lazar}{Lazar v. Superior Court}{(1996) 12 Cal.4th 631}
\caseentry{prentice}{Prentice v. North American Title}{(1963) 59 Cal.2d 618}
\caseentry{moradishalal}{Moradi-Shalal v. Fireman's Fund}{(1988) 46 Cal.3d 287}
\caseentry{actionapt}{Action Apartment Ass'n v. City of Santa Monica}{(2007) 41 Cal.4th 1232}
\caseentry{flatley}{Flatley v. Mauro}{(2006) 39 Cal.4th 299}
\caseentry{rusheen}{Rusheen v. Cohen}{(2006) 37 Cal.4th 1048}
\caseentry{shaolian}{Shaolian v. Safeco Ins. Co.}{(1999) 71 Cal.App.4th 268}
\caseentry{lucido}{Lucido v. Superior Court}{(1990) 51 Cal.3d 335}
\caseentry{crowley}{Crowley v. Katleman}{(1994) 8 Cal.4th 666}
\caseentry{mycogen}{Mycogen Corp. v. Monsanto Co.}{(2002) 28 Cal.4th 888}
\caseentry{statefarm}{State Farm Fire \& Casualty Co. v. Superior Court}{(1996) 45 Cal.App.4th 1093}
\caseentry{doctorsco}{Doctors' Company v. Superior Court}{(1989) 49 Cal.3d 39}
\caseentry{americanair}{American Airlines v. Sheppard, Mullin}{(2018) 28 Cal.App.5th 1077}
\caseentry{gdsearle}{G.D. Searle \& Co. v. Superior Court}{(1975) 49 Cal.App.3d 22}
\subsubsection*{Statutes}
\statuteentry{ccp473}{Code Civ. Proc.,}{\textsection~473}
\statuteentry{ccp436}{Code Civ. Proc.,}{\textsection~436}
\statuteentry{civ1709}{Civ. Code,}{\textsection\textsection~1709, 1710}
\statuteentry{civ3294}{Civ. Code,}{\textsection~3294}
\statuteentry{civ2338}{Civ. Code,}{\textsection~2338}
\statuteentry{bpc17200}{Bus. \& Prof. Code,}{\textsection~17200 et seq.}
\statuteentry{ins79003}{Ins. Code,}{\textsection~790.03}
\statuteentry{ccr2695}{Title 10 CCR,}{\textsection\textsection~2695.4, 2695.7}
\vspace{4mm}
\phantomsection\label{toc:i-introduction}
\section{I. INTRODUCTION}\label{i.-introduction}

Plaintiff FRANCISCUS DYLAN ROSARIO respectfully submits this Memorandum in support of his Motion for Leave to File a Third Amended Complaint (``TAC'').

\textbf{Two-Layer Cure Strategy:} This TAC motion is the \textbf{second layer} of Plaintiff's cure strategy. The \textbf{first layer} is the pending Motion for Leave to File Second Amended Complaint (filed April 16, 2026; hearing April 30, 2026). If the SAC is granted and defense brings Round D challenges, this TAC stands ready as a further cure. The proposed TAC cures the seven pleading objections raised in Defendant's Motion to Strike filed April 22, 2026, while preserving Plaintiff's meritorious claims for fraud and unfair business practices.

This case involves CSAA's fraudulent denial of a claim supported by \textbf{confession, victim report, and eyewitness corroboration}---a claim denied just 21 days after the collision, before any meaningful investigation could have been completed. The February 25, 2021 denial letter, signed by \textbf{Sarah Rash, Claims Representative}, represented that CSAA had ``concluded our investigation'' and ``carefully reviewed the facts'' when the foundational incident records had not been obtained.

California's strong policy favoring amendment and resolution on the merits requires that leave be granted.

\begin{center}\rule{0.5\linewidth}{0.5pt}\end{center}

\phantomsection\label{toc:ii-statement-of-facts}
\section{II. STATEMENT OF FACTS}\label{ii.-statement-of-facts}

\subsection{A. The Collision and the Precipitous Denial}\label{a.-the-collision-and-the-precipitous-denial}

On \textbf{February 4, 2021}, Plaintiff was involved in a motor vehicle collision in San Francisco with CSAA's insured, Subhi Abdelhalim.

On \textbf{February 25, 2021}---just 21 days after the collision---CSAA sent Plaintiff a denial letter for \textbf{Claim No.~1004-09-1054}, signed by \textbf{Sarah Rash, Claims Representative}. The letter stated:

\begin{quote}
``We have concluded our investigation of your claim. Based on our investigation, and after carefully reviewing the facts and circumstances of your claim, we have concluded that our insured is not liable\ldots{}''
\end{quote}

\subsection{B. The Three Objective Items CSAA Failed to Obtain}\label{b.-the-three-objective-items-csaa-failed-to-obtain}

At the time of the denial, CSAA had not obtained:

\begin{enumerate}
\def\labelenumi{\arabic{enumi}.}
\tightlist
\item
  \textbf{911 audio recording} containing the insured's own statements acknowledging contact;
\item
  \textbf{Plaintiff's contemporaneous 911 report} stating the driver hit him;
\item
  \textbf{Body-worn camera footage} containing independent eyewitness (Justin Rosemond) corroboration.
\end{enumerate}

Taken together, these three sources constituted virtually uncontested evidence of liability.

\subsection{C. The Three-Anchor Delayed-Discovery Sequence}\label{c.-the-three-anchor-delayed-discovery-sequence}

Plaintiff did not discover the fraud until:

\begin{enumerate}
\def\labelenumi{\arabic{enumi}.}
\tightlist
\item
  \textbf{January 17, 2025:} Defense MIL No.~17 revealed the 911 materials were of ``unknown provenance''---proving CSAA never obtained them;
\item
  \textbf{February 18, 2025:} Jeremy Jessup testified the materials were produced to defense counsel in 2022, while defense counsel stated ``this is news to me'';
\item
  \textbf{April 16, 2025:} Defense counsel admitted: ``They were provided by Dolan.''
\end{enumerate}

\begin{center}\rule{0.5\linewidth}{0.5pt}\end{center}

\phantomsection\label{toc:iii-legal-standard}
\section{III. LEGAL STANDARD}\label{iii.-legal-standard}

\subsection{A. Leave to Amend Should Be Liberally Granted}\label{a.-leave-to-amend-should-be-liberally-granted}

``The policy favoring amendment is so strong that denial of leave to amend is rarely justified.'' \emph{Howard v. County of San Diego} (2010) 184 Cal.App.4th 1422, 1428.

\subsection{B. Prejudice Is the Key Consideration}\label{b.-prejudice-is-the-key-consideration}

``The only legitimate ground for denial of leave to amend is prejudice to the opposing party.'' \emph{Kolani v. Gluska} (1998) 64 Cal.App.4th 402, 412.

\begin{center}\rule{0.5\linewidth}{0.5pt}\end{center}

\phantomsection\label{toc:iv-the-tac-cures-all-seven-defense-objections}
\section{IV. THE TAC CURES ALL SEVEN DEFENSE OBJECTIONS}\label{iv.-the-tac-cures-all-seven-defense-objections}

\subsection{Objection {\textrightarrow{}} Cure Table}\label{objection-cure-table}

{\def\LTcaptype{table} % do not increment counter
\begin{longtable}[]{@{}
  >{\raggedright\arraybackslash}p{(\linewidth - 8\tabcolsep) * \real{0.0484}}
  >{\raggedright\arraybackslash}p{(\linewidth - 8\tabcolsep) * \real{0.3065}}
  >{\raggedright\arraybackslash}p{(\linewidth - 8\tabcolsep) * \real{0.1613}}
  >{\raggedright\arraybackslash}p{(\linewidth - 8\tabcolsep) * \real{0.2419}}
  >{\raggedright\arraybackslash}p{(\linewidth - 8\tabcolsep) * \real{0.2419}}@{}}
\toprule\noalign{}
\begin{minipage}[b]{\linewidth}\raggedright
\#
\end{minipage} & \begin{minipage}[b]{\linewidth}\raggedright
Defense Objection
\end{minipage} & \begin{minipage}[b]{\linewidth}\raggedright
TAC Cure
\end{minipage} & \begin{minipage}[b]{\linewidth}\raggedright
TAC Section/¶
\end{minipage} & \begin{minipage}[b]{\linewidth}\raggedright
Key Authority
\end{minipage} \\
\midrule\noalign{}
\endhead
\bottomrule\noalign{}
\endlastfoot
1 & \emph{Trope} --- pro-per fees barred & Deletion of pro-per fees; \emph{Prentice} damages element with four line-items (incl.~\$75k cost-judgment as induced-litigation damages) & \textsection{}III, ¶¶ 22--24, 24A & \emph{Prentice} (1963) 59 Cal.2d 618, 620 \\
2 & Oppression / malice not pleaded & \emph{Lazar} particulars; deny-delay-defend; confession/victim/eyewitness triple; pleading-vs-evidentiary bright line & \textsection{}IV, ¶¶ 25--30, 25A & Civ. Code \textsection{} 3294; \emph{Lazar} (1996) 12 Cal.4th 631; \emph{G.D. Searle} (1975) 49 Cal.App.3d 22, 29 \\
3 & Fraud lacks particularity & Sarah Rash named; Claim No.~1004-09-1054; Feb.~25, 2021 date; specific representations; evidence triple with SFPD incident identifiers & \textsection{}I, ¶¶ 3--4, 8; \textsection{}IV, ¶¶ 25, 29 & \emph{Lazar} \\
4 & Litigation privilege / \emph{Ribas} & Pre-litigation carveout with three-factor \emph{Action Apartment} walk-through; \emph{Flatley} illegality carve-out; non-communicative conduct; \emph{Ribas} distinction & \textsection{}II, ¶¶ 14--21, 16A, 21A & \emph{Action Apartment} (2007) 41 Cal.4th 1232; \emph{Flatley} (2006) 39 Cal.4th 299; \emph{Rusheen} (2006) 37 Cal.4th 1048 \\
5 & Standing / \emph{Shaolian} & Direct-addressee theory; voluntary direct communication; Civ. Code \textsection{} 2338 agency/ratification spine & \textsection{}I.5, ¶¶ 5--6, 6A & Civ. Code \textsection{} 2338; \emph{Doctors' Company} (1989) 49 Cal.3d 39 \\
6 & Collateral estoppel / \textsection{} 1032 & Distinct primary right; PI disclaimer; 2021-forward anchor; \emph{Lucido} five-factor walk-through; 801/802 double-recovery reservation & \textsection{}VI, ¶¶ 41--44, 41A, 44A & \emph{Lucido} (1990) 51 Cal.3d 335; \emph{Mycogen} (2002) 28 Cal.4th 888; \emph{Crowley} (1994) 8 Cal.4th 666 \\
7 & Investigation standard not pleaded & Title 10 CCR \textsection{} 2695.7 requirements; itemized failures; Ins. Code \textsection{} 790.03(h)(5) UCL predicate & \textsection{}VII, ¶¶ 45--47, 47A & 10 CCR \textsection{}\textsection{} 2695.4, 2695.7; Ins. Code \textsection{} 790.03(h)(5); \emph{State Farm} (1996) 45 Cal.App.4th 1093 \\
8 & Anti-SLAPP / \textsection{} 425.16 overlap & Express wrongful-conduct-scope disavowal --- no litigation-conduct predicate & First CoA, ¶ 52A & CCP \textsection{} 425.16(e); \emph{Flatley} \\
9 & UCL fraudulent-prong sufficiency & Reasonable-consumer test pleaded & Third CoA, ¶¶ 66, 66A & \emph{Williams v. Gerber} (9th Cir. 2008) 552 F.3d 934; \emph{Chapman v. Skype} (2013) 220 Cal.App.4th 217 \\
\end{longtable}
}

\begin{center}\rule{0.5\linewidth}{0.5pt}\end{center}

\phantomsection\label{toc:v-specific-answers-to-defenses-legal-arguments}
\section{V. SPECIFIC ANSWERS TO DEFENSE'S LEGAL ARGUMENTS}\label{v.-specific-answers-to-defenses-legal-arguments}

\subsection{A. Standing Is Established (Shaolian Rebuttal)}\label{a.-standing-is-established-shaolian-rebuttal}

Defense cites \emph{Shaolian v. Safeco Ins. Co.} (1999) 71 Cal.App.4th 268. The TAC cures this by:

\begin{enumerate}
\def\labelenumi{\arabic{enumi}.}
\tightlist
\item
  Pleading the \textbf{direct-addressee theory} (TAC \textsection{}I.5, ¶¶ 5--6): Plaintiff is the direct addressee of a specific misrepresentation communicated by CSAA to Plaintiff individually;
\item
  Distinguishing \emph{Shaolian} (policy benefits against an insurer in a coverage dispute) from this action (common law fraud by the person to whom the representation was made);
\item
  Establishing that CSAA \textbf{voluntarily undertook a direct communicative relationship} with an unrepresented claimant.
\end{enumerate}

\subsection{B. Litigation Privilege Does Not Apply}\label{b.-litigation-privilege-does-not-apply}

Defense invokes Civil Code \textsection{} 47(b). The TAC cures this by:

\begin{enumerate}
\def\labelenumi{\arabic{enumi}.}
\tightlist
\item
  Identifying the February 25, 2021 denial as a \textbf{pre-litigation} communication (TAC \textsection{}II, ¶¶ 14--16)---no litigation was pending or contemplated;
\item
  Affirmatively walking through the three \emph{Action Apartment} contemplation factors (TAC ¶ 16A) and pleading each factor's absence as of February 25, 2021: (a) no litigation pending or contemplated; (b) no counsel retained, no demand, no tolling agreement; (c) no preservation letter and no litigation reserve;
\item
  Anchoring the \emph{Flatley v. Mauro} (2006) 39 Cal.4th 299, 317--21, illegality carve-out (TAC ¶ 21A)---\textsection{} 47(b) categorically does not protect conduct that is illegal as a matter of law, including fraud under Civ. Code \textsection{}\textsection{} 1709, 1710;
\item
  Pleading \textbf{non-communicative conduct} (failure to investigate) that falls outside the privilege (TAC \textsection{}II, ¶¶ 20--21);
\item
  Distinguishing \emph{Ribas}---which addresses IIED, not fraud (TAC \textsection{}II, ¶¶ 17--19).
\end{enumerate}

\subsection{C. Fraud Is Pleaded with Particularity}\label{c.-fraud-is-pleaded-with-particularity}

The TAC satisfies \emph{Lazar} by alleging:

\begin{itemize}
\tightlist
\item
  \textbf{WHO:} Sarah Rash, Claims Representative
\item
  \textbf{WHAT:} False representation of investigative completeness
\item
  \textbf{WHEN:} February 25, 2021 (21 days after collision)
\item
  \textbf{WHERE:} CSAA claims office; communication to Plaintiff in San Francisco
\item
  \textbf{HOW:} Written denial letter without investigation
\item
  \textbf{BY WHAT MEANS:} U.S. Mail
\end{itemize}

\subsection{D. UCL Claim Is Proper Under Moradi-Shalal Carve-Out}\label{d.-ucl-claim-is-proper-under-moradi-shalal-carve-out}

The TAC expressly disclaims any private \textsection{} 790.03 cause of action. The UCL claim is predicated on \textbf{independently actionable common law fraud} (TAC \textsection{}VII, ¶ 47). \emph{State Farm Fire \& Casualty Co.~v. Superior Court} (1996) 45 Cal.App.4th 1093, 1103--04.

\subsection{E. UCL Injunction Is Properly Targeted}\label{e.-ucl-injunction-is-properly-targeted}

The TAC narrows the UCL injunction to prohibit CSAA from ``\textbf{issuing denial letters stating an investigation is `concluded'} when foundational incident records have not been obtained or reviewed'' (TAC \textsection{}V(h), ¶ 38)---tracking the defense's own FAC-cite at p.18:5-6 and the actual \textsection{} 2695.7 violation.

\subsection{F. Corporate Agency and Ratification}\label{f.-corporate-agency-and-ratification}

The TAC's new ¶ 6A embeds a Civil Code \textsection{} 2338 agency and ratification spine that directly answers the anticipated defense argument that ``CSAA did not itself act.'' Under \textsection{} 2338, CSAA is responsible for the wrongful acts and willful omissions of its agents --- Sarah Rash, claims department personnel, and any retained claims professionals acting within the scope of the claims-handling agency --- as a matter of statutory principal-agent liability. The managing-agent / ratification framework pleaded in TAC ¶¶ 26--27 is reinforced by ¶ 6A's express allegation that CSAA ratified the February 25, 2021 denial by failing to correct it after the January 17, 2025 / February 18, 2025 / April 16, 2025 record events. See \emph{Doctors' Company v. Superior Court} (1989) 49 Cal.3d 39 (insurer directly liable when it acts outside ordinary insurer role). This closes the prior ``CSAA did not itself act'' imputation gap noted in the Cure Matrix.

\subsection{G. Primary Rights Preclusion Framework}\label{g.-primary-rights-preclusion-framework}

The TAC's new ¶ 41A walks through each factor of the California Supreme Court's issue-preclusion test in \emph{Lucido v. Superior Court} (1990) 51 Cal.3d 335, 341, against the CGC-21-594102 record, showing that (i) the issues are not identical (negligence in 21-594102 vs.~claim-handling fraud here); (ii) the claim-handling fraud issue was not actually litigated in 21-594102; (iii) the 9--3 negligence verdict did not necessarily decide any fraud issue; (iv) the 21-594102 judgment is not final because it is on appeal (A173827); and (v) CSAA was not a named party in 21-594102 and privity does not extend to independent claim-handling fraud. Under the primary-rights doctrine of \emph{Mycogen Corp.~v. Monsanto Co.} (2002) 28 Cal.4th 888, 904, and \emph{Crowley v. Katleman} (1994) 8 Cal.4th 666, 681--82, the right to be free from commercial fraud (this action) is a different primary right from the right to be free from negligent driving (21-594102). The new ¶ 44A separately reserves against double recovery between 631802 and 631801 by disclaiming overlap between legal-damages relief here and equitable-vacatur relief in 631801, with pre-judgment election of remedies preserved.

\begin{center}\rule{0.5\linewidth}{0.5pt}\end{center}

\phantomsection\label{toc:vi-no-prejudice-to-defendant}
\section{VI. NO PREJUDICE TO DEFENDANT}\label{vi.-no-prejudice-to-defendant}

The proposed amendments do not prejudice Defendant because:

\begin{enumerate}
\def\labelenumi{\arabic{enumi}.}
\tightlist
\item
  \textbf{No New Parties:} The TAC names only CSAA and DOES 1-50.
\item
  \textbf{No New Theories:} Fraud, UCL, and concealment were already alleged.
\item
  \textbf{No Additional Discovery Needed:} The TAC is based on the same factual foundation.
\item
  \textbf{No Trial Delay:} The case is at an early stage.
\end{enumerate}

\begin{center}\rule{0.5\linewidth}{0.5pt}\end{center}

\phantomsection\label{toc:vii-conclusion}
\section{VII. CONCLUSION}\label{vii.-conclusion}

For the foregoing reasons, Plaintiff respectfully requests that this Court grant leave to file the Third Amended Complaint in the form lodged herewith.

\begin{center}\rule{0.5\linewidth}{0.5pt}\end{center}

\textbf{Dated:} \_\_\_\_\_\_\_\_\_\_\_\_\_\_\_

\textbf{FRANCISCUS DYLAN ROSARIO}\\
Plaintiff, In Pro Per\\
7624 Melody Drive\\
Rohnert Park, CA 94928\\
Tel: 650.488.7510\\
Email: soltrinox@gmail.com

\nolinenumbers
% Generated by build_tac_response_802_court_pdfs.write_tac_signature_block_tex()
\vspace{0.5\baselineskip}
\begin{singlespace}
\begin{flushleft}
Dated: April 22, 2026 \\
\vspace{0.25in}
\noindent\includegraphics[height=0.5in]{../../../Support/signature.png}\\[\baselineskip]
FRANCISCUS DYLAN ROSARIO \\
Plaintiff, In Pro Per \\
\vspace{0.15in}
7624 Melody \\
Rohnert Park, CA 94928 \\
650.488.7510 \\
E: soltrinox@gmail.com
\end{flushleft}
\end{singlespace}


\end{document}
